\section{Hyperparameter tuning}\label{hptuning}
We have learned about various augmentation techniques in chapter \ref{Augmentations}. However, the possibilities of image manipulation are endless, which means that there are an unlimited number of augmentation methods to consider. And each augmentation technique may have hyperparameters that control the magnitude of the transformation, leading to even more possibilities for image manipulation. To tackle the challenge of selecting the appropriate augmentation techniques and hyperparameter values for a specific dataset, it is critical to adopt a systematic approach. This involves having prior knowledge of the dataset and a solid understanding of the augmentation techniques.

In this section, we will discuss the importance of choosing augmentation techniques as well as hyperparameter tuning in augmentation. Then, we will present a systematic approach for tuning these hyperparameters to optimize the image transformation pipeline for the specific dataset and model architecture. In chapter \ref{results}, we will represent the optimal hyperparameter values in a table format for each augmentation technique.

\subsection{Choosing the right augmentation techniques}
Unfortunately, there is no shortcut for selecting the appropriate image augmentation techniques as this process heavily depends on the characteristics of the dataset used to train the model. The optimal set of augmentations can vary significantly based on the distribution, size, and complexity of the dataset. While a similar policy may be referenced if the datasets are similar in size and object, in our case, the image size is relatively large (256x256 pixels) compared to other popular aerial image datasets, such as BigEarthNet (120x120 pixels) \cite{sumbul2019bigearthnet} or Airbus Wind Turbines Datasets (128x128 pixels) \cite{windmilldataset}. Additionally, the object of interest is small compared to the background, necessitating an empirical exploration of new augmentation policies.

To demonstrate the significance of exploring augmentation policies for different datasets, let us consider the example of CenterCrop. CenterCrop transformation can be meaningful in datasets with center biases and small objects. However, applying it to other datasets can remove important details and make training more challenging. The selection of appropriate image augmentation techniques should consider the characteristics of the dataset and the specific use case to improve model performance. Figure \ref{fig:hptuningfail_1} demonstrates the impact of applying the CenterCrop transformation inappropriately in Airbus Wind Turbines Windmill Dataset \cite{windmilldataset}.

\begin{figure}
	 \centerline{\includegraphics[width=0.8\textwidth]{src/approach/hptuning_example1.png}}
	 {\caption{Illustration of the CenterCrop Transformation  2/5 in our data and one sample from Airbus Windturbines Dataset \cite{windmilldataset}. The image on the bottom right illustrates the windmill being barely recognizable after the transformation, highlighting the importance of appropriate selection of image augmentation techniques for the specific use case.}\label{fig:hptuningfail_1}}
\end{figure}

Therefore, it is important to carefully consider the characteristics of the dataset and the specific learning task when choosing the right augmentation technique. This can involve exploring a wide range of augmentations methods and evaluating the performance of the model on a validation set to determine the optimal combination of hyperparameters.

However, due to computational resource limitations, exploring all possible augmentation techniques is not feasible. If a transformation is expected to not enhance robustness to known datasets, the technique will have lower priority to be explored. In this context, a good example for our datasets is RandomSunFlare, as illustrated in Figure \ref{fig:randomsunflare}. RandomSunFlare is an image augmentation technique that adds a bright patch or flare to the image, simulating the effect of the sun or other light source shining on the camera lens. However in our case, as the camera is often mounted on a drone or other airborne device and positioned high above the ground. Hence, the angle and position of the sun or light source relative to the camera are likely to be different from those in ground-level photography, making the addition of a sun flare unrealistic and potentially misleading. As a result, other augmentation techniques, such as rotation, translation, or color jitter, may be more promising for increasing the diversity and variability of the dataset. 

\begin{figure}
	 \centerline{\includegraphics[width=0.5\textwidth]{src/approach/RandomSunFlare.png}}
	 {\caption{Illustration of the RandomSunflare Transformation.  While we cannot be certain about its effects, we do not find the simulation of sunlight meaningful for our aerial image dataset. Therefore, this method is not a promising option and has a lower priority for exploration.}\label{fig:randomsunflare}}
\end{figure}

\subsection{Choosing the right hyperparmeters}
After considering the appropriate augmentation techniques, we also have to pick the right parameter for the augmentation. In chapter \ref{Augmentations}, we learned about the various image augmentation techniques used in this work and how they have different sets of parameters that control different aspects of the transformation, such as magnitude, change limit, filter size, threshold, and more. Although every augmentation technique incorporates a certain degree of randomness, the range of the randomness is not learned during training, but rather set manually by the user at the beginning. Therefore, the choice of hyperparameters can significantly impact the quality and diversity of the augmented images, which in turn can affect the performance of the model.

To conserve computational resources, it is important to mention that not all hyperparameter values are meaningful or relevant to explore. For example, in the Solarize method, pixels exceeding a particular threshold are inverted using the formular:
\begin{equation}
\text{max value} - \text{pixel value} \quad \text{if} \quad \text{pixel value} \geq \text{threshold}
\end{equation}
Exploring extremely small values, such as zero, is not informative as the resulting image may resemble an inverted picture, which does not provide the model with any new features, because this can be figured by the model itself by adapting the convolutional filter kernels \cite{krizhevsky2017imagenet}. 

On the other hand, it is also non-optimal to explore extreme hyperparameter values that are highly likely to introduce non-label-preserving transformations. For instance, when exploring contrast in the CLAHE method, exploring values of clipping limit greater than 100, may often result in too noisy images as in Figure \ref{fig:clahe_fail}, hence reduce the model performance. Therefore, it is essential to carefully select the hyperparameter values to explore to optimize the performance of the model while minimizing the computational cost. 

In hyperparameter tuning algorithms, it is crucial to define a search space beforehand, which sets bounds for all hyperparameters and incorporates prior knowledge of them. However, using a uniform distribution to explore values within a range could be inefficient since it assumes that all values in the range are equally important. To solve this problem, a log-uniform distribution is used, which explores values uniformly across several orders of magnitude. For instance, when searching for the best value for a hyperparameter such as learning rate, exploring a range from 0.1 to 1e-5 using a uniform distribution may result in spending more time exploring many more values in the range of 0.1 and 0.2 than in the range of 1e-4 and 2e-4, because the interval between 0.1 and 0.2 is 1000 times bigger. This can be inefficient especially for hyperparameters like the learning rate that have a logarithmic relationship with the model's performance metrics. This approach also applies to augmentation hyperparameter search, for example, when tuning the clipping limit in CLAHE. Rather than exploring values from 4 to 8, which would require seven runs, exploring values such as 4, 8, 16, and 32 helps to pinpoint the optimal range.

\begin{figure}
	 \centerline{\includegraphics[width=0.8\textwidth]{src/approach/CLAEH_fail.png}}
	 {\caption{Illustration of the CLAHE Transformation with Clipping limit as 100. This demonstrates that when applying extreme hyperparameters, the resulting augmented images can be overly noisy and the objects in the images may become unrecognizable.\label{fig:clahe_fail}}}
\end{figure}

\subsection{Methods for tuning hyperparameters}
In order to select the optimal hyperparameters for a particular image augmentation pipeline, various search methods can be used. There are many advanced techniques for hyperparamter tuning which is different in complexity and computational resource requirement. Two of the most common search methods are Grid Search and Random Search, which are both simple to implement and computationally feasible. In this section, we introduce and compare these search methods and discuss their advantages and disadvantages in hyperparameter tuning for image augmentation. 

\subsubsection{Grid Search}
The ideas of grid search is fairly simple.  We construct a grid of potential discrete hyperparameter values and train the model with every possible combination. We record the performance of the model for each set of hyperparameters and select the combination that yields the best performance. This can be accomplished using nested for-loops through predefined values. To illustrate this concept, we utilize an example with the CLAHE Transformation \ref{CLAHE}. As shown in Figure \ref{fig:clahe_hp_grid}, we explore the clipping limit and filter size values each of 2, 4, 8, 16, and 32.

\begin{figure}
	 \centerline{\includegraphics[width=0.8\textwidth]{src/approach/clahe_hp.png}}
	 {\caption{Ilustration of Grid Search in CLAHE. The size of each dot represents the F1 score for each variable combination. As illustrated, there are several obvious minima and some combinations that are relatively comparable in efficiency. The combination with the best F1 score, which is 0.119422, is selected. This combination consists of a clipping limit of 16 and a filter size of 4.}\label{fig:clahe_hp_grid}}
\end{figure}

Grid Search is an exhausive algorithm that explores all possible combinations, allowing it to find the optimal point within the search space. However, the downside is that it can be slow. Checking every combination within the space requires a significant amount of time, which may not always be available or sometimes even impossible.

\subsubsection{Random Search}
Random Search is an alternative to Grid Search that randomly samples the search space instead of discretizing it into a Cartesian grid. Unlike Grid Search, Random Search does not have a fixed endpoint but requires a predetermined time budget that determines the number of trials. Random Search has been found to be particularly effective, especially when some hyperparameters have a much larger range of variation than others, which can result in a non-cubic search space \cite{bergstra2012random}. As an example, we can consider the hyperparameter tuning of Color Jittering, which consists of four parameters: Brightness, Hue, Contrast, and Saturation. Each parameter has a range between 0.2 and 0.8.. Table \ref{tab:random_search} shows the results of 4 iterations.

\begin{table}[ht]
\centering
\caption{Hyperparameter exploration with Random Search in Color Jittering. The maximum values of brightness, hue, contrast, and saturation are uniformly selected within the range of 0.2 and 0.8. After eight runs, the chosen maximum values are: 0.8 for brightness, 0.42 for hue, 0.54 for contrast, and 0.65 for saturation.}
\begin{tabular}{lccccc}
\toprule
Index & Brightness & Hue & Contrast & Saturation & F1 Score\\
\midrule
1 & 0.57 & 0.45 & 0.56 & 0.78 & 0.0125 \\
2 & 0.74 & 0.73 & 0.6 & 0.4 & 0.0250 \\
3 & 0.41 & 0.28 & 0.55 & 0.47 & 0.0388 \\
4 & 0.55 & 0.2 & 0.23 & 0.64 & 0.0418 \\
5 & 0.67 & 0.26 & 0.65 & 0.27 & 0.0507 \\
6 & 0.77 & 0.77 & 0.24 & 0.36 & 0.0518 \\
7 & 0.48 & 0.55 & 0.34 & 0.64 & 0.0519 \\
8 & 0.8 & 0.42 & 0.54 & 0.65 & \textbf{0.0585} \\
\bottomrule
\end{tabular}
\label{tab:random_search}
\end{table}

Random search is a valuable option when dealing with several hyperparameters with a continuous values. By selecting a subset consisting of 5-50 randomly chosen points, depending on the computational resources available, we can obtain a set of hyperparameter values that are reasonably good. In our experiment, we only use 5-10 points. While this set may not necessarily be the best point, it can still provide a good model.

\subsubsection{Other tuning techniques}
In addition to grid search and random search, there are other hyperparameter tuning techniques such as Bayesian optimization, genetic algorithms, and gradient-based optimization. 

\begin{itemize}
\item Bayesian optimization \cite{frazier2018tutorial}: This advanced hyperparameter tuning approach uses Bayesian inference to construct a probabilistic model of the objective function. The objective is to create a model that accurately approximates the performance of the model for any given set of hyperparameters. The Bayesian optimization then selects the next set of hyperparameters to evaluate using this model, aiming to maximize the model's performance while minimizing the number of evaluations needed. The main advantage of Bayesian optimization is that it can quickly identify promising regions of the hyperparameter space, allowing for efficient exploration of a large number of potential hyperparameter configurations. However, it is computationally expensive and requires expertise in selecting the prior distribution and the acquisition function. 
\item Genetic algorithms \cite{orive2014evolutionary} \cite{aszemi2019hyperparameter}: This technique is inspired by natural selection and uses a population of candidate solutions that evolve over time through crossover and mutation to find the optimal set of hyperparameters. The main advantage of genetic algorithms is their ability to search a large space of potential hyperparameters quickly and efficiently. But this approach is very computationally expensive, and it tends to find multiple local optima rather than one global optimum.
\item Gradient-based optimization \cite{bengio2000gradient}: This approach involves computing the gradients of the objective function with respect to the hyperparameters and then using these gradients to iteratively update the hyperparameters until the optimal values are found. Gradient-based optimization algorithms require the objective function to be differentiable with respect to the hyperparameters, which is sometimes infeasible.
\end{itemize}

While these advanced techniques  are known to be more optimal in finding the optimal combination of hyperparameters, they tend to be computationally expensive and require a high level of expertise to design. As a result, in this thesis, we will focus solely on grid search and random search as they are simpler and more straightforward to implement.